\textbf{Proof.} Put \(r=r_b=b/(b+1)\), \(s=1/(b+1)\), \(q_b=r^{-b}=(1+1/b)^b\), and \(\lambda=r-s=(b-1)/(b+1)\). The proposed assignment process is the stationary symmetric two-state Markov chain with transition matrix
\[
\begin{pmatrix}r&s\\s&r\end{pmatrix}.
\]
It is outcome-independent and therefore belongs to \(\mathcal D_T\). Stationarity and the Markov property give
\[
\upsilon_{t,b,a}=\frac12r^{k_t}.
\]
Since \(k_t\le b\),
\[
\upsilon_{t,b,a}\ge\frac12r^b=\frac{1}{2q_b}>\frac{1}{2e}>0,
\]
where \(q_b<e\) follows from \(b\log(1+1/b)<1\). For every \(b\), equality in the weak inequality is obtained whenever \(T\ge b+1\) and \(t=b+1\). Since \(q_b\to e\), the uniform infimum is \(1/(2e)\); it is approached but never attained at a finite \(b\). Thus every denominator in \(\widehat\vartheta_b\) is positive.

Fix \(Y\in\mathcal P_{T,\beta}(C)\), write \(y_{t,a}=Y_t(\mathbf a_{1:t})\), and define \(Z_t\) and \(\widetilde\theta_b\) as in the statement. The two exposure events at a fixed time are disjoint, and the inclusion-probability identity gives
\[
\mathbb E Z_t=y_{t,1}-y_{t,0},
\qquad
\mathbb E\widetilde\theta_b=\operatorname{GATE}_T(Y).
\]
Moreover, boundedness and \(\upsilon_{t,b,a}=r^{k_t}/2\) give
\[
\mathbb E Z_t^2
=2r^{-k_t}(y_{t,1}^2+y_{t,0}^2)
\le4B^2r^{-k_t}.
\tag{1}
\]
The bound in (1) is also an exact variance bound when \(y_{t,0}=y_{t,1}=B\), because then \(\mathbb E Z_t=0\).

Let \(u=t+h\) with \(1\le h\le b\). The two homogeneous windows overlap. Opposite-arm exposure events are disjoint, whereas, for either arm \(a\), stationarity and the Markov property yield
\[
\Pr(H_{t,b,a}H_{u,b,a}=1)=\frac12r^{k_t+h}.
\]
Since \(k_u\ge h\), put \(\rho=2r^{h-k_u}\ge2\). Direct expansion, with no Cauchy--Schwarz relaxation, gives
\[
\begin{aligned}
\operatorname{Cov}(Z_t,Z_u)
={}&(\rho-1)(y_{t,1}y_{u,1}+y_{t,0}y_{u,0})\\
&+y_{t,1}y_{u,0}+y_{t,0}y_{u,1}.
\end{aligned}
\tag{2}
\]
Indeed, the joint same-arm contribution to \(\mathbb E(Z_tZ_u)\) is \(\rho(y_{t,1}y_{u,1}+y_{t,0}y_{u,0})\), while
\[
\mathbb EZ_t\,\mathbb EZ_u
=(y_{t,1}-y_{t,0})(y_{u,1}-y_{u,0}).
\]
Using \(\rho\ge2\) and \(\lvert y_{v,a}\rvert\le B\) in (2),
\[
\lvert\operatorname{Cov}(Z_t,Z_u)\rvert
\le2(\rho-1)B^2+2B^2
=4B^2r^{h-k_u}.
\tag{3}
\]
Equality in (3) holds when all four endpoint values equal \(B\). Also, since \(k_{t+h}\le k_t+h\), one has \(2h+k_t-k_{t+h}>0\); because \(0<r<1\), the exact factor in (3) is strictly smaller than the corresponding near-pair Cauchy--Schwarz factor for every nonempty near pair.

Now suppose \(u-t>b\), and put \(d=u-b-t\ge1\). Conditional on \(W_{u-b}\), the last \(b\) coordinates in the homogeneous window remain equal to that coordinate with probability \(r^b\). Hence
\[
\mathbb E(Z_u\mid W_{u-b}=1)=2y_{u,1},
\qquad
\mathbb E(Z_u\mid W_{u-b}=0)=-2y_{u,0}.
\]
The transition identity for the symmetric chain is
\[
\Pr(W_{t+d}=j\mid W_t=i)
=\frac12\{1+(2i-1)(2j-1)\lambda^d\}.
\]
It follows that
\[
\left|\mathbb E(Z_u\mid W_{1:t})-\mathbb EZ_u\right|
=\lambda^d|y_{u,1}+y_{u,0}|
\le2B\lambda^d.
\]
Since \(Z_t\) is measurable with respect to \(W_{1:t}\) and \(\mathbb E|Z_t|=|y_{t,1}|+|y_{t,0}|\le2B\),
\[
|\operatorname{Cov}(Z_t,Z_u)|\le4B^2\lambda^d.
\tag{4}
\]
When \(y_{v,0}=y_{v,1}=B\) for every \(v\), the signs in this calculation align and equality holds in (4).

There are \(T-h\) near pairs at lag \(h\), and \((T-b-d)_+\) far pairs at distance index \(d\). Adding (1), twice the covariance sums in (3)--(4), and interpreting empty sums as zero gives
\[
\operatorname{Var}\left(\sum_{t=1}^T Z_t\right)
\le B^2V_{T,b}^{\dagger}.
\tag{5}
\]
For the constant endpoint array \(y_{t,0}=y_{t,1}=B\), equality holds simultaneously in every diagonal, near-pair, and far-pair term used in (5). This endpoint array is induced by the admissible constant potential-outcome array. Therefore
\[
\sup_{Y\in\mathcal P_{T,\beta}(C)}
\operatorname{Var}_{\pi^{(b)}}\left(\sum_{t=1}^T Z_t\right)
=B^2V_{T,b}^{\dagger}.
\tag{6}
\]
As an exact smallest-horizon check, when \(T=b=1\), the constant endpoint array gives \(Z_1=2B\) or \(-2B\) with equal probabilities, so its variance is \(4B^2=B^2V_{1,1}^{\dagger}\). When \(T=2\) and \(b=1\), the four equiprobable assignment paths give \(B^{-1}\sum_tZ_t\in\{-6,-2,2,6\}\), whose variance is \(20=V_{2,1}^{\dagger}\). This exactness concerns only the auxiliary homogeneous-history variance. Unbiasedness and Jensen's inequality give the risk implication
\[
\mathbb E\left|\widetilde\theta_b-\operatorname{GATE}_T(Y)\right|
\le\frac BT\sqrt{V_{T,b}^{\dagger}}.
\tag{7}
\]

We next establish the uniform coefficient. Since \(r^{-k_t}\le q_b\),
\[
4\sum_{t=1}^T r^{-k_t}\le4q_bT.
\tag{8}
\]
For the near pairs, \(k_{t+h}\le b\), so
\[
\begin{aligned}
\sum_{h=1}^b\sum_{t=1}^{T-h}r^{h-k_{t+h}}
&\le T\sum_{h=1}^br^{h-b}\\
&=T\sum_{j=0}^{b-1}r^{-j}
=Tb(q_b-1).
\end{aligned}
\tag{9}
\]
For the far pairs,
\[
8\sum_{d=1}^{(T-b-1)_+}(T-b-d)\lambda^d
\le8(T-b)\frac{\lambda}{1-\lambda}
=4(T-b)(b-1).
\tag{10}
\]
Combining (8)--(10) proves
\[
V_{T,b}^{\dagger}
\le4q_bT+8Tb(q_b-1)+4(T-b)(b-1).
\tag{11}
\]
For \(b\ge3\), the elementary expansion
\[
q_b=\sum_{k=0}^b\frac1{k!}\prod_{j=0}^{k-1}\left(1-\frac jb\right)
\]
and the deficits in the \(k=2\) and \(k=3\) terms give
\[
e-q_b\ge\frac1b-\frac1{3b^2}.
\]
Also \(q_b\le e<11/4\). Consequently
\[
8b(e-q_b)-4(q_b-1)
\ge8-\frac8{3b}-4(e-1)>0.
\tag{12}
\]
The difference between \((8e-4)Tb\) and the right side of (11) is exactly
\[
T\{8b(e-q_b)-4(q_b-1)\}+4b(b-1),
\tag{13}
\]
which is positive by (12). At \(b=1\), direct summation gives
\[
V_{T,1}^{\dagger}=16T-12<(8e-4)T.
\]
At \(b=2\), (11) gives
\[
V_{T,2}^{\dagger}\le33T-8<(16e-8)T.
\]
Thus, for every admissible \(T,b\),
\[
V_{T,b}^{\dagger}\le4q_bT+8Tb(q_b-1)+4(T-b)(b-1)<(8e-4)Tb.
\tag{14}
\]
For fixed \(b\), divide the three exact sums defining \(V_{T,b}^{\dagger}\) by \(Tb\) and let \(T\to\infty\). The boundary terms vanish and geometric summation gives
\[
\lim_{T\to\infty}\frac{V_{T,b}^{\dagger}}{Tb}
=\frac{4q_b}{b}+8(q_b-1)+4\frac{b-1}{b}.
\]
Since \(q_b\to e\),
\[
\lim_{b\to\infty}\lim_{T\to\infty}\frac{V_{T,b}^{\dagger}}{Tb}=8e-4.
\tag{15}
\]
Equations (14)--(15) show that \(8e-4\) is the smallest coefficient in a uniform inequality \(V_{T,b}^{\dagger}\le KTb\).

It remains to compare observed outcomes with homogeneous-history outcomes. If \(t\le b+1\), then \(k_t=t-1\), so on \(H_{t,b,a}=1\) the full history through time \(t\) equals \(\mathbf a_{1:t}\); replacement error is zero. If \(t\ge b+2\), boundedness and the historical polynomial envelope give, on \(H_{t,b,a}=1\),
\[
|Y_t(W_{1:t})-y_{t,a}|
\le\min\left\{2B,C\sum_{\ell=b+1}^{t-1}(1+\ell)^{-(1+\beta)}\right\}.
\tag{16}
\]
For each arm and time,
\[
\mathbb E\left[\frac{H_{t,b,a}}{\upsilon_{t,b,a}}\right]=1.
\]
The triangle inequality and (16) therefore yield
\[
\mathbb E|\widehat\vartheta_b-\widetilde\theta_b|
\le E_{T,\beta}^{\sharp}(b).
\tag{17}
\]
Dropping the cap in (16), interchanging the two finite sums, and comparing the decreasing tail with its integral give
\[
\begin{aligned}
E_{T,\beta}^{\sharp}(b)
&\le\frac{2C}{T}\sum_{\ell=b+1}^{T-1}(T-\ell)(1+\ell)^{-(1+\beta)}\\
&\le\frac{2C}{\beta}\frac{(T-b-1)_+}{T}(b+1)^{-\beta}.
\end{aligned}
\tag{18}
\]
All sums are zero when empty, so (16)--(18) cover \(T=1\), \(b=1\), and \(b=T\). Finally, combine (7), (17), and (18) with the triangle inequality to obtain
\[
\sup_{Y\in\mathcal P_{T,\beta}(C)}
\mathbb E_{\pi^{(b)}}|\widehat\vartheta_b-\operatorname{GATE}_T(Y)|
\le U_{T,\beta}^{\dagger}(b)
\le\sqrt{8e-4}\,B\sqrt{\frac bT}
+\frac{2C}{\beta}\frac{(T-b-1)_+}{T}(b+1)^{-\beta}.
\]
The equality in (6) does not turn either inequality here into an equality: the replacement term remains an expected error bound, and no exact bias, maximal carryover error, exact estimator risk, or minimax risk is asserted. \(\square\)